% Drawing conventions for the electrostatics diagrams.
\tikzset{
  em diagram/.style={
    x=.84cm,y=.84cm,line width=.55pt,line cap=round,line join=round,
    every node/.style={font=\small,inner sep=1.6pt},
  },
  guide/.style={line width=.35pt},
  aux/.style={guide,dash pattern=on 2.5pt off 1.65pt},
  vec/.style={-{Stealth[length=3.1pt,width=2.6pt]}},
  dim/.style={guide,{Stealth[length=2.3pt,width=2.2pt]}-{Stealth[length=2.3pt,width=2.2pt]}},
  field/.style={
    guide,postaction={decorate},
    decoration={markings,mark=at position .54 with {\arrow{Stealth[length=2.5pt,width=2.2pt]}}},
  },
}

% A coordinate, a label anchor, and the label.
\newcommand{\EMpoint}[3]{\fill (#1) circle[radius=1.15pt] node[#2] {$#3$};}
\newcommand{\EMball}[3]{\EMcharge{#1}{#2}{#3}{}}

% An open circle with its sign at the center and an optional exterior label.
\newcommand{\EMcharge}[4]{%
  \begin{scope}[shift={(#1)}]
    \draw[fill=white] (0,0) circle[radius=2.65pt];
    \if\relax\detokenize{#4}\relax\else
      \draw[guide] (-1.60pt,0)--(1.60pt,0);
      \if+#4\draw[guide] (0,-1.60pt)--(0,1.60pt);\fi
    \fi
    \if\relax\detokenize{#3}\relax\else
      \node[#2,outer sep=2.65pt] at (0,0) {$#3$};
    \fi
  \end{scope}%
}

% Horizontal dimension between two x coordinates, with extension lines.
\newcommand{\EMdim}[5]{%
  \draw[guide] (#1,#4)--(#1,#3) (#2,#4)--(#2,#3);
  \draw[dim] (#1,#3)--(#2,#3) node[midway,fill=white,inner sep=1pt] {$#5$};
}

% Angle at a vertex, in degrees, with the label outside the arc.
\newcommand{\EMangle}[5]{%
  \begin{scope}[shift={(#1)}]
    \draw[guide] (#2:#4) arc[start angle=#2,end angle=#3,radius={#4}];
    \node at ({(#2+#3)/2}:{#4+.24}) {$#5$};
  \end{scope}%
}
